\documentclass[11pt,a4paper]{article}

\usepackage[margin=2.2cm]{geometry}
\usepackage{amsmath,amssymb}
\usepackage{booktabs}
\usepackage{xcolor}
\usepackage[most]{tcolorbox}
\usepackage[colorlinks=true,linkcolor=blue!55!black,citecolor=blue!55!black,
            urlcolor=blue!55!black]{hyperref}

\newcommand{\ee}{\mathrm{e}}
\newcommand{\ii}{\mathrm{i}}
\newcommand{\dalpha}{\Delta\alpha}
\newcommand{\dchi}{\Delta\chi}
\newcommand{\ours}{\textsf{ours}}
\newcommand{\theirs}{\textsf{theirs}}

\title{\bfseries Cross-check of two independent solution sets\\[2pt]
\large Challenges 20 and 39: what agrees, what differs, and what is wrong}
\author{}
\date{}

\begin{document}
\maketitle
\vspace{-2.2\baselineskip}

\begin{center}\small
Reproduce everything here with \texttt{crosscheck/compare\_solutions.py}
(output in \texttt{crosscheck/compare\_output.txt}).
\end{center}

\section*{Verdict}

\begin{tcolorbox}[colback=green!4,colframe=green!35!black,
  title=Challenge 39 --- the two answers are the \emph{same expression}]
\texttt{sympy.simplify(theirs $-$ ours)} returns exactly $0$, and the two agree
numerically to $5\times10^{-17}$ over $10\times10$ blocks of $(n,n')$ at four
parameter sets including complex $\alpha$.  \textbf{Nothing is wrong with
either.}  The $\alpha^{n+1}$-versus-$\alpha^{n}$ appearance is two algebraic
rewritings applied at once, not a disagreement --- see Sec.~\ref{sec:39}.
\end{tcolorbox}

\begin{tcolorbox}[colback=orange!5,colframe=orange!60!black,
  title=Challenge 20 --- same physics; their $\omega_t$ keeps one term ours drops]
The two $g$ \emph{formulae are algebraically identical}, and both use the
cycle-averaged ($C=\tfrac12$) convention that is the default of our
\texttt{answer.py}.  The single substantive difference is that their
$\omega_t$ retains the optical-binding correction to the torsional stiffness,
which we derived with the identical coefficient
(\texttt{challenge20.pdf}, Sec.~5) but dropped, following the source paper.
Verified to twelve digits at six $(P_0,R)$ points:
\[
  \frac{\omega_t^{\theirs}}{\omega_t^{\ours}\sqrt{1+\delta}}=1,
  \qquad
  \frac{g^{\theirs}}{g^{\ours}(C=\tfrac12)}\sqrt{1+\delta}=1,
  \qquad
  \delta=\frac{V\chi_{\parallel}}{\pi R^{3}}\bigl(\cos kR+kR\sin kR\bigr).
\]
\textbf{Their answer is the better one}: $\delta$ is the same order in the
dipole--dipole expansion as $g$ itself, so keeping $g$ while dropping $\delta$
is inconsistent (Sec.~\ref{sec:consistency}).  Our omission is a
$1$--$3\,\%$ error of consistency, not of algebra.
\end{tcolorbox}

\noindent
What was compared:
\begin{center}\small
\begin{tabular}{lll}
\toprule
& \ours & \theirs \\
\midrule
Challenge 20 & \texttt{challenges/20/answer.py} & \texttt{solutions/20/answer.py} \\
Challenge 39 & \texttt{challenges/39/answer.py} & \texttt{solutions/39/answer.py} \\
Provenance   & this session & \texttt{codex-cli/gpt-5.6-sol} (unreviewed) \\
\bottomrule
\end{tabular}
\end{center}

\section{Challenge 39: identical, written two ways}
\label{sec:39}

The two boxed answers are
\begin{align}
  \ours:\quad
  \langle n'|\hat\rho_{c,\rm ss}|n\rangle
  &=\ee^{-|\alpha|^{2}}\delta_{n0}\delta_{n'0}
   +\frac{2\,\ee^{-|\alpha|^{2}}|\alpha|^{2}\,\alpha^{n'}(\alpha^{*})^{n}}
         {\sqrt{n!\,n'!}\left[\frac{g^{2}}{2\gamma^{2}}(n-n')^{2}+n+n'+2\right]},
  \label{eq:ours39}\\[4pt]
  \theirs:\quad
  \langle n'|\hat\rho_{c,\rm ss}|n\rangle
  &=\ee^{-|\alpha|^{2}}\delta_{n'0}\delta_{n0}
   +\frac{4\gamma^{2}\ee^{-|\alpha|^{2}}\alpha^{n'+1}(\alpha^{*})^{n+1}}
         {\sqrt{n'!\,n!}\left[g^{2}(n'-n)^{2}+2\gamma^{2}(n+n'+2)\right]}.
  \label{eq:theirs39}
\end{align}
Two rewritings, applied simultaneously, map one into the other:
\begin{equation}
  \underbrace{|\alpha|^{2}\alpha^{n'}(\alpha^{*})^{n}
  =\alpha^{n'+1}(\alpha^{*})^{n+1}}_{\text{shifts \emph{both} exponents by one}}
  \qquad\text{and}\qquad
  \underbrace{\frac{2}{\frac{g^{2}}{2\gamma^{2}}D^{2}+S}
  =\frac{4\gamma^{2}}{g^{2}D^{2}+2\gamma^{2}S}}_{\text{clears }2\gamma^{2}
  \text{ from the denominator}},
\end{equation}
with $D=n-n'$ and $S=n+n'+2$.  The first explains the $n+1$ exponents; the
second supplies the $4\gamma^{2}$ that pairs with them.  Note also
$D^{2}=(n-n')^{2}=(n'-n)^{2}$, so the ordering there is immaterial.

\paragraph{Machine verification.}
\texttt{compare\_solutions.py} evaluates the two symbolically and subtracts:
\begin{center}\small
\texttt{sympy.simplify(theirs - ours), complex alpha = 0}\\
\texttt{sympy.simplify(theirs - ours), alpha > 0\ \ \ \ \ = 0}
\end{center}
(the complex case needs the single hand substitution
$|\alpha|^{2}\to\alpha\alpha^{*}$, which SymPy will not make on its own for a
symbol of unknown phase).  Numerically, over $10\times10$ blocks:

\begin{table}[h]
\centering\small
\begin{tabular}{rrrr}
\toprule
$g$ & $\gamma$ & $\alpha$ & $\max|\theirs-\ours|$ \\
\midrule
$1.0$ & $1.0$ & $1.0$ & $2.8\times10^{-17}$ \\
$10.0$ & $1.0$ & $0.8+0.6\mathrm{i}$ & $1.2\times10^{-18}$ \\
$0.2$ & $3.0$ & $1.5$ & $5.6\times10^{-17}$ \\
$2.0$ & $0.7$ & $-0.9+0.2\mathrm{i}$ & $2.8\times10^{-17}$ \\
\bottomrule
\end{tabular}
\caption{Challenge 39: agreement at machine epsilon.}
\end{table}

\noindent
Both are additionally confirmed against a QuTiP integration of the master
equation to $5.3\times10^{-11}$ (\texttt{challenges/39/qutip\_check.py}), and
both derivations run through the same intermediate kernel
$K_{NM}=2\sqrt{NM}/[\frac{g^{2}}{2\gamma^{2}}(N-M)^{2}+N+M]$ with the same
trace check.  \textbf{No action required on Challenge 39.}

\section{Challenge 20: where the difference sits}

\subsection{The $g$ formulae are the same}

Their derivation uses the cycle-averaged pair energy
$U_{12}=-\tfrac12\operatorname{Re}[\mathbf p_1^{*}\!\cdot\!\mathsf G\!\cdot\!
\mathbf p_2]$ --- i.e.\ $C=\tfrac12$ in our notation, the choice we argued for
and made the default.  Their $J$ and our $\kappa$ agree,
\begin{equation}
  J=\frac{(\dalpha)^{2}E_0^{2}}{8\pi\epsilon_0R^{3}}s_{\perp}(kR)
  =-\tfrac12\dalpha^{2}E_0^{2}\operatorname{Re}G_{\perp}
  =\kappa\big|_{C=1/2},
  \qquad
  s_{\perp}(u)=(1-u^{2})\cos u+u\sin u,
\end{equation}
and both then set $g=J/(2I\omega_t)$.  Reducing our expression with
$V=\tfrac{4\pi}{3}ab^{2}$, $I=\rho V(a^{2}+b^{2})/5$ and
$E_0^{2}=4P_0/(\pi c\epsilon_0w_0^{2})$ gives
\begin{equation}
  g=\frac{5ab^{2}(\dchi)^{2}P_0}
         {3\pi c\rho w_0^{2}R^{3}(a^{2}+b^{2})\,\omega_t}\,s_{\perp}(kR),
\end{equation}
character for character their boxed $g$, sign included.

\subsection{Their $\omega_t$ keeps the optical-binding term}

Each particle sits not in the bare tweezer field but in the tweezer field plus
the neighbour's scattered field.  They obtain the correction from the local
field,
\begin{equation}
  \frac{|E_{\rm loc,x}|^{2}}{E_0^{2}}=1+\frac{V\chi_{\parallel}}{\pi R^{3}}
  \left(\cos kR+kR\sin kR\right)\equiv1+\delta,
  \qquad
  \kappa_t=\frac{\dalpha E_0^{2}}{2}\left(1+\delta\right).
\end{equation}
We obtained the same $\delta$ by a different route --- the $\theta_i^{2}$ part
of the pair energy $U_{12}$, giving
$\delta=4\alpha_{\parallel}(\cos kR+kR\sin kR)/(4\pi\epsilon_0R^{3})$, which
equals $(4/3)ab^{2}\chi_{\parallel}s_{\parallel}/R^{3}$, identical --- and then
dropped it, quoting it as ``${\approx}2\%$ in $\omega_t^{2}$'' in
\texttt{challenge20.pdf} Sec.~5.  The two routes agreeing is a useful check in
itself: the local-field enhancement of the drive \emph{is} the physical origin
of the $\theta_i^{2}$ cross term.

\subsection{Numbers}

Benchmark parameters $a=300$~nm, $b=180$~nm, $\rho=3500$~kg\,m$^{-3}$,
$\epsilon_r=5.7$, $w_0=500$~nm, $\lambda=1064$~nm.

\begin{table}[h]
\centering\small
\begin{tabular}{rrr rrr r}
\toprule
$P_0$ [W] & $R/\lambda$ & $\delta$
& \multicolumn{3}{c}{$\omega_t/2\pi$ [MHz]}
& ratio \\
\cmidrule(lr){4-6}
& & & \theirs & \ours & $\ours\times\sqrt{1+\delta}$
& $\theirs/(\ours\sqrt{1+\delta})$ \\
\midrule
$0.6$ & $0.996$ & $+0.02193$ & $1.053110$ & $1.041750$ & $1.053110$ & $1.000000000000$ \\
$0.6$ & $0.950$ & $-0.02652$ & $1.027842$ & $1.041750$ & $1.027842$ & $1.000000000000$ \\
$0.6$ & $2.000$ & $+0.00318$ & $1.043406$ & $1.041750$ & $1.043406$ & $1.000000000000$ \\
$1.0$ & $0.996$ & $+0.02193$ & $1.359559$ & $1.344894$ & $1.359559$ & $1.000000000000$ \\
$1.0$ & $0.950$ & $-0.02652$ & $1.326938$ & $1.344894$ & $1.326938$ & $1.000000000000$ \\
$1.0$ & $2.000$ & $+0.00318$ & $1.347031$ & $1.344894$ & $1.347031$ & $1.000000000000$ \\
\bottomrule
\end{tabular}
\caption{$\omega_t$: the entire difference is the factor $\sqrt{1+\delta}$.}
\end{table}

\begin{table}[h]
\centering\small
\begin{tabular}{rr rrr r}
\toprule
$P_0$ [W] & $R/\lambda$
& \multicolumn{3}{c}{$g/2\pi$ [kHz]} & ratio \\
\cmidrule(lr){3-5}
& & \theirs & $\ours\,(C=\tfrac12)$ & $\ours\,(C=1)$
& $\frac{\theirs}{\ours}\sqrt{1+\delta}$ \\
\midrule
$0.6$ & $0.996$ & $-38.6869$ & $-39.1088$ & $-78.2176$ & $1.000000000000$ \\
$0.6$ & $0.950$ & $-41.4889$ & $-40.9350$ & $-81.8700$ & $1.000000000000$ \\
$0.6$ & $2.000$ & $-19.7620$ & $-19.7934$ & $-39.5868$ & $1.000000000000$ \\
$1.0$ & $0.996$ & $-49.9446$ & $-50.4892$ & $-100.9785$ & $1.000000000000$ \\
$1.0$ & $0.950$ & $-53.5619$ & $-52.8468$ & $-105.6937$ & $1.000000000000$ \\
$1.0$ & $2.000$ & $-25.5126$ & $-25.5531$ & $-51.1063$ & $1.000000000000$ \\
\bottomrule
\end{tabular}
\caption{$g$: identical formulae; the residual $\lesssim1\%$ is the same
$\sqrt{1+\delta}$ entering through $\omega_t$ in the denominator.  The last
column shows the source paper's convention $C=1$ for reference --- a factor of
two away from \emph{both} of us.}
\label{tab:g20}
\end{table}

\subsection{Why their choice is the consistent one}
\label{sec:consistency}

Expanding the pair energy to quadratic order in the tilts,
\begin{equation}
  U_{12}=\text{const}
  +\underbrace{\tfrac12\operatorname{Re}[G_{xx}]E_0^{2}\alpha_{\parallel}\dalpha
    \left(\theta_1^{2}+\theta_2^{2}\right)}_{\textstyle\to\ \delta}
  -\underbrace{\tfrac12\operatorname{Re}[G_{yy}]\dalpha^{2}E_0^{2}
    \theta_1\theta_2}_{\textstyle\to\ g},
\end{equation}
both terms are \emph{first order} in the dipole--dipole interaction and
$O(V^{2}/R^{3})$: they enter at the same order.  Numerically at
$R=1.06\,\mu$m the frequency shift is $\delta/2=1.10\%$ while
$g/\omega_t=3.75\%$ --- comparable.  Keeping the second while discarding the
first, as we did (following arXiv:2504.08194), is therefore not a controlled
approximation.  The next omitted order, $O(V^{3}/R^{6})$ multiple scattering,
is genuinely smaller and is dropped by both of us.

\section{So what is actually wrong?}

\paragraph{In our solution.}
\begin{enumerate}\itemsep2pt
\item \textbf{$\omega_t$ is incomplete.}  Multiply our $\omega_t$ by
$\sqrt{1+\delta}$ with $\delta=(V\chi_{\parallel}/\pi R^{3})(\cos kR+kR\sin
kR)$, and recompute $g$ with that $\omega_t$.  We derived $\delta$ and quoted
its size correctly; we should not have dropped it.
\item \textbf{Interface mismatch.}  \texttt{solutions/20/answer.py} matches the
official \texttt{template.py} signature
\texttt{answer(a, b, rho, k, epsilon\_r, P\_0, w\_0, R)} with $a,b,w_0,R$ in
nm and $P_0$ in mW; ours uses SI keyword arguments and a different signature.
For submission ours needs a thin wrapper.  Challenge 39 has the same issue:
the template wants a single SymPy expression from
\texttt{answer(n, np, g, gamma, alpha)}; ours returns NumPy matrices.
\end{enumerate}

\paragraph{In their solution.}
No algebraic or physical error was found.  Three things are worth knowing:
\begin{enumerate}\itemsep2pt
\item \textbf{It does not flag the factor-of-two convention.}  Both of us use
the cycle-averaged $C=\tfrac12$; the source paper arXiv:2504.08194 uses $C=1$
throughout Eqs.~(4)--(9) while cycle-averaging its trap term, so its $g_0$ is
twice ours (last column of Table~\ref{tab:g20}).  If the reference answer was
generated from that paper, both of our answers are a factor of two below it.
That is a grading risk, not a physics error --- on the physics $C=\tfrac12$ is
correct.
\item \textbf{Only one libration plane is treated.}  Their solution librates in
the $x$--$y$ plane.  Because the polarisation is parallel to $\mathbf R$ here,
the $x$--$z$ plane is exactly degenerate and carries the identical coupling, so
the stated $H$ is one of two identical copies.  This does not change $\omega_t$
or $g$ (see \texttt{challenges/20/generalized problem/}), but it is worth
stating.
\item \textbf{Unreviewed.}  \texttt{metadata.json} records
\texttt{review\_status: generated-unreviewed}.
\end{enumerate}

\paragraph{Three candidate answers, ranked.}
\begin{center}\small
\begin{tabular}{llll}
\toprule
& $C$ & binding term in $\omega_t$ & $g/2\pi$ at $1$~W, $R=1.06\,\mu$m \\
\midrule
Source paper arXiv:2504.08194 & $1$ & no & $-100.98$~kHz (they quote $119$) \\
\ours\ as shipped & $\tfrac12$ & no & $-50.49$~kHz \\
\theirs\ (most complete) & $\tfrac12$ & yes & $-49.94$~kHz \\
\bottomrule
\end{tabular}
\end{center}

\noindent
On physics the ranking is \theirs\ $>$ \ours\ $>$ paper.  Against an automated
grader keyed to the paper it would be the reverse, which is why our
\texttt{answer.py} exposes the \texttt{convention} switch; the binding term
should be added to it as well.

\end{document}
